\documentclass[14pt, russian]{article}
\usepackage{graphicx}
\linespread{1.3}
\usepackage{amsbsy}
\usepackage[T2A]{fontenc}
\usepackage[cp1251]{inputenc}
\usepackage[russian]{babel}
\usepackage{epstopdf}
\usepackage[autostyle]{csquotes}
\usepackage[a4paper,left=2.5cm,right=2.5cm,top=2.5cm,bottom=2.5cm]{geometry}%
\usepackage{cuted}   
\usepackage{lipsum}  
\usepackage[intoc,nocfg,russian]{nomencl}
\newcommand{\nomencl}[2]{#1 --- #2\nomenclature{#1}{#2}}
\setlength{\nomlabelwidth}{3em}
\setlength{\nomitemsep}{-\parsep}
\setcounter{tocdepth}{2}
\renewcommand{\nomlabel}[1]{#1 ---}
\makenomenclature
\usepackage{wrapfig}
\let\vec=\mathbf
\def\mprp{\mbox{\tiny $\bot$}}
\def\mprl{\mbox{\tiny $\|$}}
\def\th{\mbox{th}}
\def\D{\mathrm{d}}
\def\beq{\begin{eqnarray}}
\def\eeq{\end{eqnarray}}
\def\ee{\varepsilon}
\def\lm{\lambda}
\newcommand{\ii}{\mathrm{i}} % мнимая единица
\newcommand{\dd}{\mathrm{d}} % дифференциал
\newcommand{\eee}{\mathrm{e}} % основание натуральных логарифмов

\newcommand{\bs}{\boldsymbol}

\begin{document}
\begin{center}
Ответ на повторный отзыв рецензента на работу\\
РЕЗОНАНСНОЕ КОМПТОНОВСКОЕ РАССЕЯНИЕ В ЗАМАГНИЧЕННОЙ СРЕДЕ\\
авторов Д. А. Румянцев, А. А. Ярков, Д. М. Шленев.

\vspace{5mm}
Уважаемая редакция
\end{center}

Благодарим рецензента за ценные замечания. В ответ на них нами были внесены в 
статью следующие исправления:
\begin{enumerate}
	\item В приложении 1 было явно указано, что в приближении узкого пика 
	используется только первое слагаемое в разложении (46)
	\item Для большей наглядности было использовано другое обозначение ${\cal 
	S}$-матричного элемента во избежание путаницы с обозначением для 
	пропагатора.
	\item Исправлены указанные опечатки в приложении 1.
\end{enumerate}

\vspace{5mm}
С уважением,

Д. А. Румянцев, А. А. Ярков, Д. М. Шленев.
%
%Функция
%%
%\begin{equation}
%\label{f}
%f_{\varepsilon} (x) = \frac{1}{\pi} \frac{\varepsilon}{x^2 + \varepsilon^2} \; 
%(\varepsilon >0)
%\end{equation}
%%
%образует дельта-образную последовательность при $\varepsilon\to 0$ [1]':
%%
%\begin{equation}
%\lim_{\varepsilon\to 0} \frac{1}{\pi} \frac{\varepsilon}{x^2 + \varepsilon^2} 
%= \delta(x).
%\end{equation}
%%
%Поскольку в формуле (21) данной статьи ширина резонанса имеет конечное 
%значение, определяемое параметром $\varepsilon = E_n''\Gamma^{s''}_n/2$, 
%квадрат модуля пропагатора может быть преобразован к виду, подобному 
%выражению~(\ref{f}), но это подобие понимается в смысле обобщенных функций, т. 
%е. в смысле сходимости к дельта-функции при интегрировании:
%\begin{eqnarray}
%	I(\varepsilon) = \pi \varphi(0) + \frac{\varepsilon}{4} P.V. 
%\int_{-\infty}^{\infty} \frac{\varphi(x)}{x^2} dx- \frac{\pi\varepsilon^2}{2} 
%\left.\frac{\partial^2 \varphi(x)}{\partial x^2}\right|_{x=0} + 
%O(\varepsilon^3),
%\end{eqnarray}
%где интеграл вычисляется в смысле главного значения (P.V. -- Principal Value), 
%$\Phi(s)$ -- Фурье-образ функции $\varphi(x)$.
%
%Таким образом замена
%\begin{eqnarray}
%	&&\left|\frac{1}{(p+q)^2_{\mprl}-M_n^2-\ii E_n''\Gamma^{s''}_n/2}\right|^2 
%= 
%	\frac{1}{((p+q)^2_{\mprl}-M_n^2)^2 + (E_n''\Gamma^{s''}_n/2)^2} = 
%	\\
%	\nonumber
%	&&=\frac{1}{E_n''\Gamma^{s''}_n/2}
%	\frac{E_n''\Gamma^{s''}_n/2}{(p+q)^2_{\mprl}-M_n^2-\ii 
%E_n''\Gamma^{s''}_n/2} 
%	\simeq\frac{2\pi}{E_n''\Gamma_n^{s''}} \delta((p+q)^2_{\mprl}-M_n^2)
%\end{eqnarray}
%с точностью до $E_n''\Gamma^{s''}_n\ll \left|(p+q)^2_{\mprl}-M_n^2\right|$ 
%имеет смысл только в подынтегральном выражении. 
%
%Реальная часть в точке резонанса равна нулю $(p+q)^2_{\mprl}-M_n^2=0$, что 
%соответствует тому, что виртуальный электрон приобретает массу $M_n$ и 
%становится реальным.
%
%В связи с этим было добавлено приложение 1, поясняющее выражение (3), 
%используя преобразование Фурье, а также уточнён текст, который указывает на 
%последующее интегрирование формулы (21), что и позволяет сделать эффективную 
%замену на дельта-функцию.
%
%1. Гельфанд, И.М. Обобщенные функции и действия над ними/ И. М. Гельфанд, Г. 
%Е. Шилов. – Москва.:Государственное издательство физико-математической 
%литературы: 1958. – 470 с.
%


\end{document}